\centerline{\bf \Huge pqr-метод}

\begin{flushright}
        \scriptsize{}
\end{flushright}

\textbf{Очень важно!} Первые три задачи безумно важны для решения методом $pqr$. Важно \textbf{САМИМ} вывести эти формулы, запомнить их и уметь быстро восстанавливать без шпаргалки.

\problem{1}
Пусть $p = a+b$, $q = ab$. Выразите через $p$ и $q$ выражения $a^2+b^2, a^3+b^3, a^4+b^4,(a-b)^2$.


\problem{2}
Для трёх переменных $a, b, c$ введём обозначения $p=a+b+c, q=a b+b c+c a$, $r=a b c$ - элементарные симметрические многочлены от трёх переменных. Выразите через $p, q, r$ выражения $a^2+b^2+c^2, a^2 b+a^2 c+b^2 a+b^2 c+c^2 a+c^2 b$, $a^3+b^3+c^3,(a b)^2+(b c)^2+(c a)^2, a^4+b^4+c^4,(a+b)(b+c)(c+a)$.

\problem{3}
а) Для целых неотрицательных $k$ положим $s_k=a^k+b^k+c^k$. Выведите рекуррентную формулу (с коэффициентами, зависящими от $p, q, r$ ), выражающую $s_k(k \geqslant 3)$ через $s_{k-1}, s_{k-2}$ и $s_{k-3}$.

\noindent
б$^{**}$) Выведите из предыдущего пункта, что всякий симметрический многочлен от трёх переменных является многочленом от элементарных симметрических

\textit{Следующие три задачи вы уже решали в предыдущем блоке в листочке с неравенствами. Теперь предлагается их перерешать с помощью $pqr$-метода сравнить сложность решений.}

\problem{4}
Докажите, что при $x, y, z \in[0,1]$ выполнено неравенство $x^2+y^2+z^2 \leqslant x y z+2$.

\problem{5}
Положительные числа $a, b, c$ удовлетворяют соотношению $a b c=1$. Докажите, что $a^2+b^2+c^2 \geqslant a+b+c$.


\problem{6}
Для положительных $a, b, c$ верно, что $a+b+c=a b+b c+c a$, докажите, что

$$
a+b+c+1 \geqslant 4 a b c .
$$

\problem{7}
$\bigl[\textbf{Russia, TST 2015}\bigr]$
Известно, что $a, b, c \geqslant 0$ и $1+a+b+c=2 a b c$. Докажите, что

$$
\frac{a b}{1+a+b}+\frac{b c}{1+b+c}+\frac{c a}{1+c+a} \geqslant \frac{3}{2} .
$$

\newpage

\hspace{3mm}

\problem{8}
$\bigl[\textbf{Mathematik-Olympiade in Deutschland, Bundesrunde, Olympiadeklasse 12}\bigr]$
Man beweise, dass für alle nicht negativen Zahlen $x, y, z$ mit $x+y+z=1$

$$
1 \leq \frac{x}{1-y z}+\frac{y}{1-z x}+\frac{z}{1-x y} \leq \frac{9}{8}
$$

gilt.

\problem{9}
$\bigl[\textbf{Iran, TST, 1996}\bigr]$
Неотрицательные числа $a, b, c$ таковы, что никакие два из них не равны нулю. Докажите, что

$$
(a b+b c+c a)\left(\frac{1}{(a+b)^2}+\frac{1}{(b+c)^2}+\frac{1}{(c+a)^2}\right) \geqslant \frac{9}{4} .
$$